\documentclass{article}
\usepackage{amsmath}
\usepackage{amssymb}
\usepackage{amsthm}
\usepackage{physics}


\newtheorem*{conjecture}{Conjecture}
\newtheorem*{theorem}{Theorem}
\newtheorem*{lemma}{Lemma}


\makeatletter
\newcommand*{\rom}[1]{\expandafter\@slowromancap\romannumeral #1@}
\makeatother


\begin{document}


\section*{Infinite summations}

\subsection*{2.1 The Archimedean Understanding}

\subsubsection*{Base length plot}
To the figure 2.1, remember that base is \(2\), not \(1\). You add the coordinates \((-1, 0), (1, 0)\).
So it's just \(|-1| + |1| = 2\).


\subsubsection*{Method of exhaustion}
The method of calculating areas by summing inscribed triangles is often referred to as
the “method of exhaustion.


\subsubsection*{Definition: Archimedean understanding of an infinite series}
Definition of his:
\begin{equation*}
    \forall L, M \in \mathbb{R}, \ \exists n \in \mathbb{N}, \ \forall m \geq n, L < S_{m} < M
\end{equation*}

\par Almost a modern one, the modern one would be:

\begin{equation*}
    \forall\varepsilon>0,~\exists n\in \mathbb{N},~\forall m\geq n:~|S_m - T|<\varepsilon
\end{equation*}

so \(L = T - \epsilon\) and \(M = T + \epsilon\).

\end{document}